\section{Why yet another seam}
\label{sec:whyseam}

\code{Interface} made \code{Stub} and \code{Spi} interchangeable, and thereby \code{EchoNode}
testable. But \code{Spi} itself lies \emph{below} that seam. Without something more it can only be
run on real hardware.

Work out what that would mean. \code{Spi} is finished in \kapref{ch:driver}. The hardware exists
in \kapref{ch:bringup}. In between you would have an untried class that nobody has run a single
time --- and the first time it is run, it is run on a board another group is debugging at the same
time, over four wires that may or may not be in the right place.

That is not a good way to find a bit-shifting error.

\begin{cppcode}
namespace driver::can
{
/**
 * @brief Moves single bytes to and from the CAN controller, and frames transactions.
 *
 *        Knows nothing about registers, command bytes or CAN: it only carries bytes.
 */
class ByteTransport
{
public:
    virtual ~ByteTransport() noexcept = default;

    virtual void select() noexcept   = 0;   /**< Assert SS; begin a transaction. */
    virtual void deselect() noexcept = 0;   /**< Release SS; end the transaction. */

    [[nodiscard]] virtual std::uint8_t transfer(std::uint8_t byte) noexcept = 0;
};
} // namespace driver::can
\end{cppcode}

Three methods. That is the whole layer.

\textbf{Why \code{transfer()} returns a byte.} SPI is full duplex; a signature like
\code{void transfer(std::uint8_t)} would lie about what the hardware does.

\textbf{Why \code{select()}/\code{deselect()} are methods of their own.} The \code{SS} framing is
part of the \emph{protocol}, not of how a byte is transferred. The transport should not have to
know that a transaction is five bytes long. Keep it dumb, and it does one thing only --- and is
easy to implement correctly both on the target and in a test.

\textbf{Why \code{noexcept} everywhere.} The layer runs on a freestanding target without
exceptions. Writing that in the signature already now makes it impossible to sneak in something
that throws.

\section{The constants}
\label{sec:constants}

\begin{cppcode}
class Spi final : public Interface
{
    // ...
private:
    /** Register indices, as sent in the command byte. */
    static constexpr std::uint8_t RegStatus{0U};
    static constexpr std::uint8_t RegTxId{1U};
    static constexpr std::uint8_t RegTxDlc{2U};
    static constexpr std::uint8_t RegTxDataLo{3U};
    static constexpr std::uint8_t RegTxDataHi{4U};
    static constexpr std::uint8_t RegTxSend{5U};
    static constexpr std::uint8_t RegRxId{6U};
    static constexpr std::uint8_t RegRxDlc{7U};
    static constexpr std::uint8_t RegRxDataLo{8U};
    static constexpr std::uint8_t RegRxDataHi{9U};
    static constexpr std::uint8_t RegRxAck{10U};
    static constexpr std::uint8_t RegErrorFlags{11U};
    static constexpr std::uint8_t RegTxAbort{12U};

    /** STATUS register bit masks. */
    static constexpr std::uint32_t StatusTxReady{1U << 0U};
    static constexpr std::uint32_t StatusRxValid{1U << 1U};
    static constexpr std::uint32_t StatusError{1U << 2U};

    /** Command byte: bit 7 set means write. */
    static constexpr std::uint8_t CmdWrite{0x80U};

    /** Value written to the write-triggered registers. */
    static constexpr std::uint32_t Trigger{0x1U};

    ByteTransport& myTransport; /**< Transport used for every register access. Not owned. */
};
\end{cppcode}

\textbf{Indices, not offsets.} The offset column of the register map is a legacy of the
memory-mapped variant. What goes out on the wire is the index, that is, $\text{offset}/4$. Write
the constants as indices directly, and you avoid a division on every call and a class of bugs
where somebody divides twice or not at all.

\textbf{\code{1U << 0U} and not \code{0x01U}.} The mask is written as a shift because the register
map is written in bit positions. \code{StatusRxValid{1U << 1U}} can be compared line by line
against the table; \code{0x02U} cannot.

\section{The two helper functions}
\label{sec:helpers}

\begin{cppcode}
void Spi::writeReg(const std::uint8_t index, const std::uint32_t value) const noexcept
{
    myTransport.select();

    static_cast<void>(myTransport.transfer(static_cast<std::uint8_t>(CmdWrite | index)));
    static_cast<void>(myTransport.transfer(static_cast<std::uint8_t>(value >> 24U)));
    static_cast<void>(myTransport.transfer(static_cast<std::uint8_t>(value >> 16U)));
    static_cast<void>(myTransport.transfer(static_cast<std::uint8_t>(value >> 8U)));
    static_cast<void>(myTransport.transfer(static_cast<std::uint8_t>(value)));

    myTransport.deselect();
}

std::uint32_t Spi::readReg(const std::uint8_t index) const noexcept
{
    myTransport.select();

    static_cast<void>(myTransport.transfer(index));

    std::uint32_t value{};
    value |= static_cast<std::uint32_t>(myTransport.transfer(0U)) << 24U;
    value |= static_cast<std::uint32_t>(myTransport.transfer(0U)) << 16U;
    value |= static_cast<std::uint32_t>(myTransport.transfer(0U)) << 8U;
    value |= static_cast<std::uint32_t>(myTransport.transfer(0U));

    myTransport.deselect();
    return value;
}
\end{cppcode}

Read \code{writeReg()} against \code{81 00 00 01 23} in \kapref{ch:spi}. Four things to note:

\textbf{Every \code{static_cast<void>} is deliberate.} \code{transfer()} is \code{[[nodiscard]]},
and during a write the return value is meaningless. The cast tells both the compiler and the
reader that this is known.

\textbf{In \code{readReg()} the first byte is discarded.} The slave loads its response shifter at
the \emph{end} of the command byte. Four \code{transfer()} calls contribute to the value, not five.

\textbf{The command byte is just \code{index} on a read.} Bit 7 low means a read; no
\code{CmdWrite}.

\textbf{The methods are \code{const}.} They do not change the state of the \code{Spi} object ---
the reference itself does not change. That the hardware changes is another matter; \code{const} in
C++ speaks about the object, not about the world. That is why \code{hasError() const} can call
\code{readReg()}.

\subsection{The sign extension bug}

This is the most common bug in the project, and it deserves a heading of its own:

\begin{cppcode}
value |= myTransport.transfer(0U) << 24U;                             // wrong
value |= static_cast<std::uint32_t>(myTransport.transfer(0U)) << 24U; // right
\end{cppcode}

\code{transfer()} returns a \code{std::uint8_t}. In the first line it is promoted to \code{int} ---
32 bits \emph{with sign} --- and shifted 24 steps. For a value with bit 7 set, that writes into the
sign bit. The result is undefined behaviour, and in practice a negative number that is then
converted to \code{std::uint32_t} with \code{0xFF} bytes where there should be none.

\textbf{The symptom: the driver works for small values and breaks as soon as a byte has bit 7
set.} So it works for \code{0x123} and breaks for \code{0xAABB}.

\begin{aside}
\note[Rule] Cast \textbf{before} the shift, not after. The same rule applies to the packing in
\kapref{ch:driver}.
\end{aside}

\subsection{Why the helper functions are private}

They are implementation details. Nobody outside the class should be able to poke at the registers,
for the same reason that \code{Interface} has no \code{readStatusRegister()}: the day somebody
does it in the application code, the architecture is broken.

The supplied test suite tests them anyway --- through \code{send()} and \code{receive()}, and by
inspecting which bytes came out. Testing a private method directly is almost always a sign that it
ought to be public, or that the test is testing the wrong thing.

\subsection{What the layer does not do}

\code{readReg()} and \code{writeReg()} are pure translations of the protocol and contain no logic
at all: no validation, no status check, no notion of what a register means.

All the meaning lives in the layer above, and that is \kapref{ch:driver}. The moment
\code{writeReg()} starts containing "if the index is \code{RegTxSend}, first check that..." two
layers have melted together, and both became harder to test.
